%! TEX root = main.tex  % Specifies the main document file for multi-file projects

\chapter{Introduction to JUCE}

JUCE is a cross-platform C++ framework designed for audio applications and plugin development. Projucer is an integrated project management tool that simplifies JUCE-based development.

Key features of Projucer:

\begin{itemize}
    \item Supports multiple operating systems (Windows, macOS, Linux, iOS, Android)
    \item Generates project files compatible with popular IDEs (Xcode, Visual Studio, Makefile)
    \item Manages project settings, source files, and dependencies
    \item Provides a live build feature for real-time code compilation and preview
\end{itemize}

Projucer organizes projects using predefined templates, which can be further modified in an external IDE.

\section{Workflow of Application Development with JUCE}

The general workflow for developing an application with JUCE includes:

\begin{enumerate}
    \item Setting up the project using Projucer
    \item Including the required JUCE and external libraries
    \item Developing, building, and deploying the project in the preferred IDE
\end{enumerate}

When using the live build functionality, the build and development process can be completed within Projucer.

\section{Project Types in Projucer}

Projucer supports several types of projects, each suited for different applications:

\begin{itemize}
    \item GUI Application: A minimal JUCE project with an empty application window, used for graphical user interface development.
    \item Animated Application: Similar to a GUI application but optimized for graphical animations.
    \item OpenGL Application: Adds OpenGL support for rendering 3D graphics and shaders.
    \item Console Application: A command-line application without a graphical interface.
    \item Audio Application: A project pre-configured for audio input and output handling.
    \item Audio Plug-In: A template for VST, AudioUnit, and AAX plugin development.
    \item Static Library: Used to create reusable software libraries linked statically.
    \item Dynamic Library: Similar to a static library but linked dynamically.
\end{itemize}

\section{Example Projects in Projucer}

JUCE provides example projects located in the \texttt{JUCE/examples} directory. These projects serve as templates and can be accessed via:

\begin{supercollider*}
File > Open Example
\end{supercollider*}

\section{Audio Plugin Development with JUCE}

\subsection{Core Plugin Classes}

A JUCE audio plugin project is built around two main classes:

\begin{itemize}
    \item \texttt{nameOfTheProjectAudioProcessor}: Inherits from \texttt{AudioProcessor} and handles audio processing and MIDI input/output.
    \item \texttt{nameOfTheProjectAudioProcessorEditor}: Inherits from \texttt{AudioProcessorEditor} and manages the plugin’s graphical user interface.
\end{itemize}

These classes implement a communication mechanism between the Digital Audio Workstation (DAW) and the plugin.

\subsection{Key Methods in the AudioProcessor Class}

The \texttt{AudioProcessor} class defines the core functionality of the plugin:

\begin{itemize}
    \item \texttt{nameOfTheProjectAudioProcessor()}: Constructor
    \item \texttt{\textasciitilde nameOfTheProjectAudioProcessor()}: Destructor
    \item \texttt{void prepareToPlay(double sampleRate, int samplesPerBlock) override}: Initializes the processor
    \item \texttt{void processBlock(AudioBuffer<float>\&, MidiBuffer\&) override}: Handles audio and MIDI processing
    \item \texttt{void releaseResources() override}: Cleans up resources when the plugin is deactivated
\end{itemize}

\subsection{Key Methods in the AudioProcessorEditor Class}

The \texttt{AudioProcessorEditor} class is responsible for rendering the GUI:

\begin{itemize}
    \item \texttt{nameOfTheProjectAudioProcessorEditor(nameOfTheProjectAudioProcessor\&)}: Constructor
    \item \texttt{\textasciitilde nameOfTheProjectAudioProcessorEditor()}: Destructor
    \item \texttt{void paint(Graphics\&) override}: Renders visual elements
    \item \texttt{void resized() override}: Adjusts the layout when the window is resized
\end{itemize}

\section{Creating GUI Controls in the Plugin}

To add a new GUI control (e.g., a slider), the following steps must be followed:

\begin{enumerate}
    \item Declare the control as a private member in \texttt{AudioProcessorEditor}.
    \item Set the control properties.
    \item Use \texttt{addAndMakeVisible} to display the control in the plugin window.
    \item Adjust the control’s size and position in the \texttt{resized()} function.
    \item Attach a listener to update the controlled parameter when the control value changes.
\end{enumerate}

Example methods:

\begin{itemize}
    \item \texttt{AudioProcessorEditor::AudioProcessorEditor()}: Initializes the control
    \item \texttt{AudioProcessorEditor::resized()}: Adjusts control layout
    \item \texttt{AudioProcessorEditor::typeOfComponentValueChanged()}: Handles control interactions
\end{itemize}

Common GUI components include:

\begin{itemize}
    \item Button
    \item TextButton
    \item Slider (including rotary sliders)
    \item ComboBox
\end{itemize}
